\documentclass{scrartcl}
\usepackage[utf8]{inputenc}
\usepackage{fancyhdr}
\usepackage{graphicx}
\usepackage{dirtytalk}
\usepackage[portrait, margin=0.5in, top=0.4in, bottom=0.3in]{geometry}
\usepackage{amsmath, amssymb}
\setlength{\parskip}{1em}
\setlength{\parindent}{0cm} 
\usepackage[mdthm,simplethm]{test}
\usepackage[utf8]{inputenc}
\usepackage{float}
\usepackage{physics}
\setlength{\parindent}{0pt}

\title{Lecture 2: Analytic Geometry and Lines}
\author{albert ye}
\date{May 2022}

\begin{document}
\maketitle
\section{Outline}
\begin{enumerate}
	\item Motivation by picture
	\item DEfinition using equations
		\begin{enumerate}
			\item Examples
			\item Classification
		\end{enumerate}
	\item Linear Transformation: Quadratic forms
\end{enumerate}

\section{Conic Sections}
As slope of the plane increases, we get ellipses, parabolas, and hyperbolas. These are all
cross-sections of a cone, or in the case of a hyperbola two cones arranged in hourclass shape.

A line in $\RR_2$ can be generalized as $a_1x_1 + a_2x_2 = a_0$, and similarly, a plane in $\RR^3$ 
can be generalized as $a_1x_1 + a_2x_2 + x_3x_3 = a_0$. A cone is given by $x_1^2 + x_2^2 = x_3^2$. 
So a conic section is defined by all $(x_1, x_2, x_3)$ s.t. $a_1x_1 + a_2x_2 + a_3x_3 = 0$ and
$x_1^2 + x_2^2 = x_3^2$.

\begin{definition}{Conic Curve}
	A \vocab{conic curve} in $\RR^2$ is $\{(x_1, x_2) \in \RR^2 | Ax_1^2 
		+ Bx_1x_2 + Cx_2^2 + Dx_1 + Ex_2 + F = 0$.
\end{definition}

Basic equations for each curve:
\begin{enumerate}
	\item Hyperbola: $x_1^2 - x_2^2 = 1$
	\item Ellipse: $x_1^2 + x_2^2 = 1$
	\item Parabola: $x_1^2 = x_2$
\end{enumerate}

\section{Linear Transformation}

\begin{definition}
	A \vocab{linear transformation} $T : \RR^2 \to \RR^2$ uses a matrix to turn
a vector $\begin{bmatrix}x_1 \\ x_2 \end{bmatrix}$ into $\begin{bmatrix}a_{11} & a_{12}
\\ a_{21} & a_{22} \end{bmatrix} \begin{bmatrix}x_1 \\ x_2\end{bmatrix}$. 
\end{definition}

\begin{proposition}
	A linear transformation $T : \RR^2 \to \RR^2$ is determined by how it acts on the basis vectors.
\end{proposition}

\begin{proof}
	Any vector $\vec v \in \RR^2$ can be written as $\vec v = a_1\vec e_1 + a_2\vec e_2$ by definition
	of basis. By linearity of $T$, we have $T(\vec v) = T(a_1 \vec e_1 + a_2 \vec e_2) = 
	T(a_1 \vec e_1) + T(a_2 \vec e_2) = a_1 T(\vec e_1) + a_2 T(\vec e_2)$. Hence, $T(\vec v)$ is 
	determined.
\end{proof}

A linear transformation $T$ transforms the bases into $T(e_1) = a_{11} \vec e_1 + a_{12} \vec e_2$
and $T(e_2) = a_{21} \vec e_1 + a_{22} \vec e_2$.

In general: $AB_{i, j} = \sum A_{i, k} B_{k, j}$.

\end{document}
